\chapter{Datasets}
\label{ch:datasets}
% Opening paragraph: four datasets, two synthetic (ours), two from the
% literature. Only LOFAR has real telescope data with human labels -- say up
% front that this is the one that matters most.

\section{Synthetic dataset}
\label{sec:synthetic}
% Write about your generator. Sources: PART 8, PART 9, "Verified facts"
% section of the context doc, Synthetic Dataset*/dataset_statistics.txt.

\subsection{Generation}
Real radio data with reliable pixel-level RFI labels is scarce, so the first
part of this work used a synthetic dataset that I generated myself. The
generator is written in Python and uses NumPy's random number generator with a
fixed seed (42), so the whole dataset can be regenerated bit for bit.

Each image is a spectrogram: one axis is frequency and the other is time, and
each pixel holds the power received in that frequency channel at that moment.
The background is not plain random noise. It is built from a sky pedestal that
varies across frequency, a smooth astronomical signal, a slow drift in time,
and Gaussian noise on top. RFI is then injected into this background. Six
types of RFI are modelled, chosen to cover the shapes commonly seen in real
observations (\cref{tab:rfi_types}). Each image receives a random selection of
types at random strengths, and images are spread over four density levels:
clean (8.2\%), light (30.0\%), moderate (36.2\%) and heavy (25.6\%).

\begin{table}[H]
  \centering
  \caption{The six RFI types in the synthetic dataset, with the number of the
    1000 images (1024$\times$265 version) that contain each type.}
  \label{tab:rfi_types}
  \begin{tabular}{llr}
    \toprule
    Type & Appearance in the spectrogram & Images \\
    \midrule
    Narrowband      & thin line at one frequency, lasting for some time  & 889 \\
    Scattered       & small isolated hits scattered over the image        & 899 \\
    Blob            & patchy, cloud-like region                           & 673 \\
    Persistent band & group of channels affected for the whole observation & 588 \\
    Broadband block & rectangle covering many channels for a while        & 488 \\
    Wideband burst  & short flash across many channels at once            & 310 \\
    \bottomrule
  \end{tabular}
\end{table}

The dataset contains 1000 images, split into 700 for training, 150 for
validation and 150 for testing. I checked that no image appears in more than
one split by comparing SHA-256 hashes of the raw pixels; there are no shared
hashes.

\subsection{Deciding which pixels are RFI}
\label{sec:synthetic_label}
Because the RFI is injected by the generator, its exact strength at every
pixel is known. A pixel is labelled as RFI when the injected RFI power there
is greater than 0.5 times the standard deviation of the local noise. Pixels
where RFI was added but is weaker than this are labelled clean, because at
that level there is effectively nothing left in the data for a detector to
find.

This rule gives labels that are noise-free and unambiguous: every labelled
pixel is, by construction, measurably above the noise. In the 276$\times$600
version, 90.9\% of the RFI pixels lie at or above one noise standard
deviation. This makes the task considerably easier than labelling real data,
a point that becomes important in \cref{ch:drivers}.

The rule was refined between the two versions of the dataset. In the
276$\times$600 version a pixel was labelled RFI if it lay inside the geometric
outline of an injected RFI shape \emph{or} exceeded the 0.5$\sigma$ threshold.
Because this was a union, it could mark pixels that carried almost no RFI
power (0.06\% of labelled pixels in the training split, and up to 4.49\% in
the worst single image). In the 1024$\times$265 version only the threshold is
used, and it is applied \emph{after} the instrument response
(\cref{sec:bandpass}) and against the total local noise of each channel, so
RFI that the instrument has genuinely buried below the noise is not
labelled.

\subsection{Two versions of the dataset}
\label{sec:synthetic_versions}
Two versions of the dataset were used (\cref{tab:synth_versions}). My first
images were 1024$\times$1024 pixels. At that size only one image fitted in a
training batch on the 6\,GB GPU, which made training unstable. Akeret et
al.~\cite{akeret2017} trained \tfunet{} on images of 276$\times$600 pixels, so
I regenerated the dataset at exactly that size to reproduce their setup as
closely as possible. When the image size changes, the RFI shapes have to be
rescaled as well: a blob 40 channels wide covers 3.9\% of a 1024-channel band
but 14.5\% of a 276-channel band. Seven size-dependent parameters of the
generator were therefore scaled in proportion to the image dimensions.

% TODO (ask Ronak): why 1024x265 for the second version? Reason not
% recorded in the repo. Add one or two sentences here.
The second version, 1024$\times$265, adds a model of the receiver's frequency
response, described in the next section.

\begin{table}[H]
  \centering
  \caption{The two versions of the synthetic dataset. Both use seed 42 and a
    700/150/150 split.}
  \label{tab:synth_versions}
  \begin{tabular}{lccc}
    \toprule
    Version & Size (freq $\times$ time) & Instrument bandpass & RFI pixels \\
    \midrule
    v3 & 276$\times$600  & no  & 14.67\% mean (test split 15.08\%) \\
    v4 & 1024$\times$265 & yes & 14.19\% mean \\
    \bottomrule
  \end{tabular}
\end{table}

\subsection{The instrument bandpass (v4)}
\label{sec:bandpass}
% - Real receivers do not amplify all frequencies equally. Formula
%     B(f) = max( W(f) * [1 + R(f) + S(f)], B_floor )
%   W = Tukey edge taper, R = slow ripple, S = standing wave.
%   [PART 9; code dataset_v4_bandpass/generate_dataset_v4.py]
% - spectrogram = B(f) * (sky + RFI) + receiver noise. Noise is added AFTER
%   the gain, as in a real receiver chain.
% - HONESTY: each ingredient is standard and citable, but the combined
%   formula is ours. Say so. Cite Harris 1978 (Tukey window), Popping & Braun
%   2008 (standing waves), Hamaker et al. 1996 (measurement equation).
%   [PART 9 PROVENANCE]

\section{HERA dataset}
\label{sec:hera}
% - Simulated HERA data from Mesarcik et al., Zenodo 6724065.
%   420 train + 140 test, 512x512, boolean masks, 2.75% RFI. [PART 2]
% - FINDING: the published split leaks. Verified by SHA-256 over raw pixels:
%   382/420 unique train, 135/140 unique test, 31 test images identical to
%   training images = 22% of the test set. [PART 2]
% - After deduplication: 351 train / 135 test / 0 overlap. [PART 3]
% - Mention briefly that it turned out too easy to be informative (models
%   at 0.999 F1) [PART 3] -- but only if you discuss HERA results at all.

\section{LOFAR dataset}
\label{sec:lofar}
% Sources: PART 10, PART 11. Figures: analysis/lofar_analysis/fig_lofar_overview.png
% and fig_lofar_profiles.png already exist.

\subsection{Origin and structure}
% - Real LOFAR observations, Mesarcik et al. 2022, Zenodo 6724065. 9.3 GB.
% - Stored as a Python list of 4 arrays:
%     train images (7500,512,512,1) float32 | train masks (AOFlagger) bool
%     test images  (109,512,512,1)  float32 | test masks (human expert) bool
% - 512x512 is a random crop from 599x616.
% - Rows = time, columns = frequency. You measured this three ways
%   [PART 11.2]: 65-74% of RFI pixels sit in 16 of 512 columns.
%   (Opposite orientation to our synthetic data.)
% - Values are RAW powers, 0 to 1.4e11 -- max is ~39,000x the 99th
%   percentile. That is why plt.matshow shows a flat square. [PART 11.1]

\subsection{Audit}
% Table of what you found [PART 11.3-11.6]:
%  - Imbalance: train 1.27% RFI (1:78), test 0.766% (1:129).
%  - 35 training images are 100% flagged (dead baselines).
%  - All 109 test images are byte-identical to 109 training images.
%    Verified by MD5 and element-wise equality, largest pixel difference 0.
%    Only the labels differ. [PART 11.5; analysis/lofar_analysis/prove_test_set_leakage.py]
%  - Clean training set: 7500 - 109 - 35 = 7356; minus 735 for validation
%    = 6621 used for training.
%  - Training labels vs human labels on the same pixels: F1 0.5698403,
%    matching the paper's published 0.5698 to four decimals. This confirms
%    the data, labels and metric are the same as theirs. [PART 11.6]

\subsection{Why it is harder than the synthetic data}
% - Training labels are ~44% wrong against the human expert.
% - RFI brightness tiers, using each image's own clean pixels as the ruler
%   [PART 15; analysis/ablation_faint_recall.py lines 82-100]:
%     bright 50.1% | moderate 18.9% | faint 11.0% | invisible 20.0%
%   "Invisible" = dimmer than the median clean pixel. No brightness
%   evidence at all -- only context can find it.
% - Explain the bucketing method briefly (you understood it from the
%   step-by-step explanation).
% - Practical note: the pickle does not fit in RAM alongside other work;
%   converted once to memory-mapped arrays (28.7 MB resident). Mention in a
%   sentence, details in the appendix.
